\documentclass[12pt,a4paper]{article}
\usepackage[T1]{fontenc}
\usepackage[utf8]{inputenc}
\usepackage[english]{babel}
\usepackage{amsmath,amssymb,amsthm}
\usepackage{geometry}
\geometry{margin=2.5cm}
\usepackage{graphicx}
\usepackage{hyperref}
\hypersetup{colorlinks=true,linkcolor=blue,citecolor=blue,urlcolor=blue}
\usepackage{booktabs}
\usepackage{enumitem}

\newtheorem{theorem}{Theorem}
\newtheorem{lemma}{Lemma}
\newtheorem{proposition}{Proposition}
\newtheorem{remark}{Remark}
\newtheorem{definition}{Definition}

\title{\textbf{A Gauge Theory of Line Defects in an Elastic Continuum:\\
Full Derivations.\\
{\large (A Geometrodynamic Model of Particles as Elastic Defects)}}}
\author{A.~Yu.~Drozdov}
\date{24 October 2026}

\begin{document}
\maketitle

\begin{abstract}
\noindent
This is an extended version of the preprint on a geometrodynamic model of
particles as defects of an elastic aether. Unlike the introductory exposition,
here \emph{every} structural statement is accompanied by a full derivation:
dimensional analysis, intermediate formulas, and, where appropriate, a proof in
the form of a theorem. The aim is to make all results reproducible and
checkable by the reader ``with a pencil''. The limits of the model (what is
derived, what is assumed, what is open) are traced through each proof, not only
in a final section.
\end{abstract}

\tableofcontents
\newpage

\section{Notation and conventions}
The three-dimensional elastic continuum (the ``aether'') is described by a
displacement field $u_i(\mathbf{x},t)$. The dimensional basis is $(L,T,M)$.
Greek indices $a,b=1,2,3$ are internal (translational) indices of the gauge
group $T(3)$; Latin $i,j,k$ are spatial. Summation over repeated indices is
implied.

\begin{definition}[Basic elastic quantities]
Shear modulus $\mu$, $[ML^{-1}T^{-2}]$; density $\rho$, $[ML^{-3}]$;
transverse-wave speed $c=\sqrt{\mu/\rho}$, $[LT^{-1}]$; Burgers vector $b$,
$[L]$.
\end{definition}

\section{Gauge structure and the Maxwell dictionary}

\subsection{Distortion field and the dislocation tensor}
Plastic incompatibility is introduced by a gauge field $\phi^a_i$. The
dislocation-density tensor (Cartan torsion) is
\begin{equation}
T^a_{ij} = \partial_i \phi^a_j - \partial_j \phi^a_i .
\label{eq:torsion}
\end{equation}
The integral over a surface spanning a contour $\partial S$ gives the Burgers
vector:
\begin{equation}
b^a = \oint_{\partial S}\phi^a_i\,dx^i = \int_S T^a_{ij}\,dS^{ij}.
\label{eq:burgers}
\end{equation}

\begin{lemma}[Dimension of the magnetic charge]
\label{lem:qm}
In the elastic basis the magnetic charge, defined as the flux of the
dislocation density $q_m=\int_S T^a\,dS$, has the dimension of length.
\end{lemma}
\begin{proof}
From \eqref{eq:torsion}, $[T^a_{ij}]=[\partial\phi]=[\phi]/L$. Since $\phi^a_i$
is dimensionless (a displacement gradient, $[\partial u]=1$), we have
$[T^a]=L^{-1}$. Then $[q_m]=[T^a]\cdot[dS]=L^{-1}\cdot L^{2}=L$. $\qed$
\end{proof}

\begin{remark}
Lemma~\ref{lem:qm} removes the apparent conflict with SI, where the magnetic
charge is measured in A$\cdot$m: in a self-contained elastic theory the ampere
is not a primary unit, and the dimension of the charge is dictated by the
theory itself.
\end{remark}

\subsection{Quasi-Maxwell equations}
Introducing the dislocation current $I^a_i=\partial_t\phi^a_i-\partial_i v^a$
(where $v^a$ is the velocity), one obtains Bianchi identities structurally
identical to Maxwell's equations \cite{lazar2010}:
\begin{align}
\partial_i {}^{*}T^{a\,i} &= 0, & \text{(no sources of }T^a\text{)}\\
\epsilon_{ijk}\partial_j I^a_k + \partial_t T^a_i &= 0. & \text{(Faraday analogue)}
\end{align}
Hence the dictionary: $\mathbf{B}\leftrightarrow T^a$,
$\mathbf{E}\leftrightarrow I^a$.

\subsection{Full derivation of the quasi-Maxwell system}
We give the derivation in detail. The Lagrangian density of the linear dynamic
gauge theory of dislocations in the $T(3)$ gauge splits into elastic and
dislocation parts \cite{lazar2010}:
\begin{equation}
\mathcal{L} = \tfrac12\rho\,v^a v_a - \tfrac12 C_{ijkl}\beta^{i}_{a}\beta^{j a}
 + \tfrac12 D_{ab}\,I^a_i I^{b\,i} - \tfrac14 A_{ab}\,T^a_{ij}T^{b\,ij},
\label{eq:lagr}
\end{equation}
where $\beta$ is the elastic distortion, $C_{ijkl}$ the elasticity tensor, and
$D_{ab}$, $A_{ab}$ material response tensors (analogues of the dielectric and
magnetic permeabilities). The field strengths are the dislocation density $T^a$
\eqref{eq:torsion} and the current $I^a$; the responses are
\begin{equation}
H_a^{\ ij} = -\frac{\partial\mathcal{L}}{\partial T^a_{ij}} = A_{ab}T^{b\,ij},
\qquad
D_a^{\ i} = \frac{\partial\mathcal{L}}{\partial I^a_i} = D_{ab}I^{b\,i}.
\label{eq:constit}
\end{equation}
Varying the action $\delta\!\int\mathcal{L}\,d^4x=0$ with respect to the gauge
potential $\phi^a$ gives the equation of motion (Amp\`ere--Oersted analogue):
\begin{equation}
\partial_j H_a^{\ ij} - \partial_t D_a^{\ i} = J_a^{\ i},
\label{eq:ampere}
\end{equation}
where $J_a$ is the source current (dislocation momentum density). The two
Bianchi identities for $T^a=d\phi^a$ (closedness of the form, $d^2=0$) give the
second pair:
\begin{equation}
\partial_{[i}T^a_{jk]}=0,\qquad \partial_t T^a_{ij}+\partial_{[i}I^a_{j]}=0.
\label{eq:bianchi}
\end{equation}
Comparing \eqref{eq:constit}--\eqref{eq:bianchi} with Maxwell's equations gives
the dictionary of Table~\ref{tab:corr} (below): the strengths $T^a,I^a$
correspond to $\mathbf B,\mathbf E$; the responses $H_a,D_a$ to
$\mathbf H,\mathbf D$; and the Peach--Koehler force on a dislocation element
\begin{equation}
f_a^{\ i} = (e_a\!\lrcorner\,I^b)\,p_b^{\ i} + (e_a\!\lrcorner\,T^b)\wedge\Sigma_b
\label{eq:pk}
\end{equation}
plays the role of the Lorentz force $q\mathbf E + q\mathbf v\times\mathbf B$.
This is the formal basis of the whole analogy: the electrodynamics of the
particle-defect is the dynamics of its dislocation field.

\begin{table}[h]
\centering
\caption{Correspondence Maxwell~$\leftrightarrow$~dislocation theory (after Lazar \cite{lazar2010}).}
\label{tab:corr}
\begin{tabular}{ll}
\toprule
Electrodynamics & Dislocation theory \\
\midrule
Magnetic field $\mathbf{B}$ & dislocation density $T^a$ \\
Electric field $\mathbf{E}$ & dislocation current $I^a$ \\
Magnetic excitation $\mathbf{H}$ & pseudomoment stress $H_a$ \\
Magnetic flux $\oint\mathbf{B}\,d\mathbf{S}$ & Burgers vector $\oint T^a = b$ \\
Lorentz force & Peach--Koehler force \\
\bottomrule
\end{tabular}
\end{table}

\section{The particle: a precessing dislocation dipole}

\subsection{Kinematics of the special point}
A ring of radius $a$ lies in a fixed plane; the special point (a localized
charge $q_m=\Delta b$) runs along it: azimuth $\varphi(t)=\varphi_0+\omega t$.
Position, velocity and acceleration:
\begin{align}
\mathbf{r}(t) &= a(\cos\varphi,\ \sin\varphi,\ 0),\\
\mathbf{v}(t) &= a\omega(-\sin\varphi,\ \cos\varphi,\ 0),\quad |\mathbf{v}|=a\omega,\\
\mathbf{W}(t) &= -a\omega^2(\cos\varphi,\ \sin\varphi,\ 0),\quad |\mathbf{W}|=a\omega^2.
\label{eq:accel}
\end{align}

\subsection{Mass as elastic self-interaction}
\begin{proposition}[Dimensional closure of mass]
The self-energy of a dislocation ring $E=\tfrac12\mu b^2 a\,\ln(8a/r_c)$
(Hirth--Lothe \cite{hirth1982}), after dropping the logarithm, yields the mass
$m=\mu b^2 a/c^2$ with the correct dimension $[M]$.
\end{proposition}
\begin{proof}
$[\mu b^2 a]=(ML^{-1}T^{-2})(L^2)(L)=ML^{2}T^{-2}=[\text{energy}]$. Division by
$c^2$ $[L^2T^{-2}]$ gives $[M]$. $\qed$
\end{proof}

\subsection{Dipole from a pair of rings (antiphase)}
Two rings at $z=\pm d/2$, special points in antiphase
($\varphi_{\text{low}}=\varphi_{\text{up}}+\pi$), charges $\pm q_m$. The
instantaneous dipole moment $\mathbf{p}=q_m\mathbf{r}_+ - q_m\mathbf{r}_-$.
A direct computation:
\begin{equation}
p_z = q_m\cdot d = \text{const},\qquad
|\mathbf{p}_{xy}| = 2 q_m a\ \ (\text{rotates with }\omega).
\end{equation}
The axial component is constant (a static dipole), the transverse one precesses.

\subsection{Stability: the centrifugal-barrier theorem}
\begin{theorem}[Equilibrium of the pair]
\label{th:stab}
Let the pair be subject to attraction $U_{\text{att}}=-A/d$ (opposite charges)
and let the angular momentum $J=m a^2\omega$ be conserved. Then the effective
energy $E(d)=-A/d+J^2/(2md^2)$ has a unique minimum at
$d_{\text{eq}}=J^2/(mA)$, and as $\omega\to0$ (i.e.\ $J\to0$) the equilibrium
vanishes ($d_{\text{eq}}\to0$).
\end{theorem}
\begin{proof}
$E'(d)=A/d^2-J^2/(md^3)=0 \Rightarrow d_{\text{eq}}=J^2/(mA)$. The second
derivative $E''(d_{\text{eq}})=-2A/d^3+3J^2/(md^4)$; substituting $d_{\text{eq}}$
gives $E''=A^4 m^3/J^6>0$ --- a minimum. As $J\to0$ the barrier term vanishes,
only attraction remains, and the minimum shifts to $d\to0$ (collapse). $\qed$
\end{proof}

\begin{remark}[Annihilation]
Collapse upon loss of precession is Frenkel recombination: an interstitial loop
($+\Delta b$) and a vacancy loop ($-\Delta b$) give a total Burgers vector
$\mathbf{b}_++\mathbf{b}_-=0$ --- the perfect lattice. Thus precession (spin) is
the mechanism protecting the particle--antiparticle pair from self-annihilation.
\end{remark}

\subsection{Spin and statistics: the configuration period}
\begin{proposition}[Period of the pair and statistics]
For a pair of ordinary rings the configuration returns at
$\varphi\to\varphi+2\pi$ (a boson). For a pair of M\"obius rings, where the sign
of the Burgers vector is modulated as $\cos(\varphi/2)$, the sign is restored
only at $\varphi\to\varphi+4\pi$ (a fermion, spin $1/2$).
\end{proposition}
\begin{proof}
The twist $\cos(\varphi/2)$ is antiperiodic under $2\pi$:
$\cos((\varphi+2\pi)/2)=\cos(\varphi/2+\pi)=-\cos(\varphi/2)$, and periodic
under $4\pi$. Hence a spinor (double-valued) structure arises precisely for the
M\"obius ring. $\qed$
\end{proof}

\begin{lemma}[Frank vector of the $\cos(\varphi/2)$ twist]
\label{lem:franck}
The Burgers twist $b(\varphi)=\cos(\varphi/2)$ is equivalent to a Frank
rotation by the fixed angle $\Omega=\pi$ per circuit $\varphi:0\to2\pi$, i.e.\ a
wedge \emph{half-disclination}. The Frank index $\nu=\Omega/2\pi=1/2$, giving
spin $1/2$.
\end{lemma}
\begin{proof}
The twist phase is $\varphi_{\text{tw}}(\varphi)=\varphi/2$. The total rotation
of the Burgers direction per circuit is
$\Omega=\oint d\varphi_{\text{tw}}=\int_0^{2\pi}\tfrac12\,d\varphi=\pi$. Over
$4\pi$ the rotation is $2\pi$ (full return), reproducing the spinor
double-valuedness of the period proposition. Through the Osipov spin connection
(\texttt{arXiv:1001.4450}) $\Omega_\varphi=i\omega\sigma_3$ with flux
$\oint\Omega_\varphi\,d\varphi=2\pi\nu$ we get $\nu=1/2$ and the magnetic moment
$\mu=\nu=1/2$ as the flux of the spin connection (not a loop current). $\qed$
\end{proof}

\begin{remark}[M\"obius dipole and the link to the $SO(3)$ sector]
The electron in the model is a pair of M\"obius rings with \emph{opposite}
directions of the Burgers vector: one ring carries the jump $+\Delta b$
(pole~$N$), the other $-\Delta b$ (pole~$S$); together they give a static
magnetic dipole $q_m d$. This distinguishes the electron (M\"obius pair, period
$4\pi$, fermion) from a boson (pair of ordinary rings, period $2\pi$).
Crucially, the Burgers twist $b(\varphi)=\cos(\varphi/2)$ is a Frank rotation by
$\pi$ per circuit --- a \emph{disclination} (rotational, $SO(3)$) degree of
freedom. Thus the M\"obius sector already lies in $SO(3)$: the magnetic moment
$\mu=\nu=n/2$ is derived as the flux of the spin connection
$\oint\omega_\varphi\,d\varphi=2\pi\nu$ (Osipov, \texttt{arXiv:1001.4450}),
not as a loop current. Hydrodynamically the same object is a vortex: the
vorticity $\nabla\times v$ corresponds to the Frank density, the circulation to
the disclination strength. The precession of the special point ($\omega>0$),
required for the stability of the pair, is that vortex; a particle without
rotation collapses.
\end{remark}

\section{The magnetic moment: why a dipole, not a current}

\begin{theorem}[Vanishing of the current axial moment]
\label{th:mu0}
For a pair of rings in antiphase with opposite charges, the period-averaged
axial ``loop-current'' magnetic moment $\boldsymbol{\mu}_{\text{cur}}=
\tfrac12\oint q\,\mathbf{r}\times\mathbf{v}$ is zero.
\end{theorem}
\begin{proof}
Both rings run in the same sense, so $\mathbf{r}\times\mathbf{v}$ is the same
for both and points along $+z$: $(\mathbf{r}\times\mathbf{v})_z=a^2\omega$. But
the charges are opposite: $q_{\text{up}}=+q_m$, $q_{\text{low}}=-q_m$. The
contributions $\tfrac12 q_m a^2\omega$ and $-\tfrac12 q_m a^2\omega$ cancel:
$\mu_{\text{cur},z}=0$. $\qed$
\end{proof}

\begin{remark}
Theorem~\ref{th:mu0} means that the electron's magnetic moment in the model is
a \emph{static} dipole moment $\mu_{\text{mag}}=q_m d$, not a loop-current
moment. This is consistent with the ``current'' along the ring being magnetic,
so that applying the electric loop-current formula to it is incorrect.
\end{remark}

\begin{proposition}[Dimensional failure of the $g$-factor from the ED formula]
Attempting to define $g$ via $\mu_{\text{mag}}=g(q_m/2m)S$ in the elastic basis
gives $g$ the dimension $T/L$, not a dimensionless quantity; hence the
electrodynamic formula does not transfer literally, and the classical $g=2$ is
not obtained in the ring model (one gets $g=1$).
\end{proposition}
\begin{proof}
$[\mu_{\text{mag}}]=[q_m][d]=L^2$. $[S]=[m a^2\omega]=M L^2 T^{-1}$.
$[(q_m/m)S]=(L/M)(ML^2T^{-1})=L^3T^{-1}$. Then
$[g]=[\mu_{\text{mag}}]/[(q_m/m)S]=L^2/(L^3T^{-1})=T/L\neq 1$. $\qed$
\end{proof}

\section{The Fermi film and gravity: full derivation}

\subsection{Setup}
The Fermi film is the three-dimensional boundary $\partial\Omega$ of the
four-dimensional action integral, stretched orthogonally to the world tube. Its
contribution to the action:
\begin{equation}
S_{\partial} = \gamma\int_{\partial\Omega} d\Sigma,
\end{equation}
where $\gamma$ is the film elastic modulus $[ML^{-1}T^{-2}]=[\text{Pa}]$.

\subsection{Source of deformation}
\begin{proposition}[Positivity of the source]
The period-averaged source of film deformation from a precessing charge
\begin{equation}
\sigma = \mu b^2 \frac{\langle W^2\rangle}{c^2}
= \mu b^2 \frac{\omega^4 a^2}{c^2}
\end{equation}
is strictly positive and independent of the sign of the charge.
\end{proposition}
\begin{proof}
From \eqref{eq:accel}, $|\mathbf{W}|=a\omega^2$, hence $\langle W^2\rangle=
a^2\omega^4>0$ (average of a constant). The sign of the charge enters through
$b^2$, i.e.\ drops out. $\qed$
\end{proof}

\subsection{The inverse-square law}
\begin{theorem}[$1/R^2$ as a consequence of the Green function]
\label{th:invsq}
If the film deformation obeys the equilibrium equation of an elastic membrane
\begin{equation}
\gamma\,\nabla^2\Delta(\mathbf{r})=-\sigma\,\delta^3(\mathbf{r}-\mathbf{r}_0),
\label{eq:poisson}
\end{equation}
then the interaction potential of two sources decays as $1/R$, and the force as
$1/R^2$.
\end{theorem}
\begin{proof}
Equation \eqref{eq:poisson} is the Poisson equation. Its fundamental solution
(the Laplacian Green function in $\mathbb{R}^3$) is $G(\mathbf{r})=1/(4\pi|\mathbf{r}|)$,
whence
\begin{equation}
\Delta(\mathbf{r})=\frac{\sigma}{4\pi\gamma\,|\mathbf{r}-\mathbf{r}_0|}.
\end{equation}
The energy of a second source $\sigma_2$ in the field of the first:
$U(R)=\sigma_2\Delta(R)=\sigma_1\sigma_2/(4\pi\gamma R)\propto R^{-1}$.
The force $F=-dU/dR=\sigma_1\sigma_2/(4\pi\gamma R^2)\propto R^{-2}$.
The exponent $-2$ is fixed by the dimension of space ($n=3$: $G\sim R^{-(n-2)}$),
not by choice. $\qed$
\end{proof}

\subsection{Two stiffness scales}
\begin{proposition}[Impossibility of a single stiffness]
The modulus reproducing the electron mass and the Planck modulus of gravity are
incompatible: their ratio is $\sim 10^{84}$.
\end{proposition}
\begin{proof}
From $m_e c^2=\mu_{\text{aeth}}\,r_e^2\lambda_C$ (with $b=r_e$, $a=\lambda_C$):
$\mu_{\text{aeth}}=m_ec^2/(r_e^2\lambda_C)\approx2.7\times10^{28}$~Pa.
The Planck modulus $\gamma_P=c^7/(8\pi\hbar G^2)\approx1.8\times10^{112}$~Pa.
The ratio $\gamma_P/\mu_{\text{aeth}}\approx6.9\times10^{83}$. $\qed$
\end{proof}

\begin{remark}[Physical meaning of the moduli]
The quantity $\mu_{\text{aeth}}$ is the \emph{shear modulus of the vacuum}: the
stiffness of the aether with respect to shear deformation, of pressure
dimension $[\text{Pa}]$. It is through it that, in the $(L,T,M)$ basis (without
the ampere), the transverse-wave speed $c=\sqrt{\mu_{\text{aeth}}/\rho}$, the
particle mass $m=\mu_{\text{aeth}}b^2a/c^2$ and the Coulomb force
$F=\mu_{\text{aeth}}b^2/R\cdot(\text{geometry})$ are defined. Physically it
governs shear deformations (dislocations), from which charge and mass are
built, and sets the physics at the particle scale (of order the Compton length
$\lambda_C$). Its value $2.7\times10^{28}$~Pa is obtained by \emph{calibration}
to the electron (substituting $m_e$ and $r_e$ into the self-energy of the
dislocation ring), not derived from first principles.

The quantity $\gamma_P=c^7/(8\pi\hbar G^2)$ is the \emph{Planck stiffness of the
vacuum}, a second elastic modulus of the same dimension $[\text{Pa}]$,
responsible for the gravitational response (deformation of the Fermi film) at
the Planck scale $\ell_P=\sqrt{\hbar G/c^3}$. The gravitational constant $G$
enters it as an \emph{input} parameter. The colossal ratio
$\gamma_P/\mu_{\text{aeth}}\sim10^{84}$ is a direct expression of the weakness
of gravity and is the main unsolved problem of the model (the origin of the
hierarchy of the two scales).
\end{remark}

\begin{remark}
Thus the gravitational constant $G$ entering $\gamma_P$ is an \emph{input}
parameter of the aether stiffness. We do not claim to derive it from $(\hbar,c)$;
only the \emph{structure} is proved (Theorem~\ref{th:invsq}), not the numerical
value. This is a fundamental difference from the fabricated versions, in which
$G$ was hidden inside the stiffness and ``extracted'' back (a tautology).
\end{remark}

\begin{remark}[Physical meaning of $\gamma_P$ and the nature of the hierarchy]
The Planck modulus is the Planck \emph{energy density} divided by $8\pi$:
\begin{equation}
\gamma_P=\frac{c^7}{8\pi\hbar G^2}=\frac{1}{8\pi}\,\frac{m_P c^2}{\ell_P^3},
\qquad \ell_P=\sqrt{\hbar G/c^3},\quad m_P=\sqrt{\hbar c/G}
\end{equation}
(numerical agreement $<1\%$). This is not an arbitrary combination but the
stiffness of the vacuum at the Planck scale; the factor $8\pi$ is the same as
in the Einstein equations. Likewise $\mu_{\text{aeth}}=m_e c^2/(r_e^2\lambda_C)$
is the energy density at the electron scale. The ratio of the two moduli
\begin{equation}
\frac{\gamma_P}{\mu_{\text{aeth}}}\sim\frac{m_P}{m_e}\ \text{(with a volume factor)}\sim10^{84}
\end{equation}
reduces to the ratio of the Planck and Compton scales, i.e.\ to the smallness
of $m_e/m_P$. Thus the hierarchy of the two stiffnesses is \emph{equivalent} to
the standard \emph{hierarchy problem} ($m_e\ll m_P$): the model does not solve
it, but neither does it introduce a new arbitrary parameter --- it inherits the
well-known smallness of $m_e/m_P$.
\end{remark}

\section{Six channels: a single table with derived laws}
\label{sec:channels}
Each decay law is a consequence of the proved structure:
\begin{itemize}[nosep]
\item \textbf{Coulomb} $1/R^2$: product of Burgers charges, the direct Coulomb
kernel (Theorem~\ref{th:invsq} for the $T^a$ field).
\item \textbf{Magnetic} $1/R^4$: interaction of two static dipoles $q_m d$
(derivative of the dipole--dipole potential $\sim1/R^3$).
\item \textbf{Gravity} $1/R^2$: Theorem~\ref{th:invsq} for the source $\sigma>0$.
\item \textbf{Exchange} $1/R^3$: phase dependence $\cos2\Delta\varphi$, short-range.
\item \textbf{Weak} $e^{-R/R_W}$: massive mediator, twist chirality (parity violation).
\item \textbf{Strong} $\propto R$: fractional Burgers vector $b/3$ and a stacking-fault
ribbon (linear potential, confinement).
\end{itemize}
Two dimensionless numbers close the system:
\begin{equation}
\alpha=\frac{ke^2}{\hbar c}=\frac{1}{137.04},\qquad
\alpha_G=\Big(\frac{m_e}{m_{\text{Planck}}}\Big)^2=1.75\times10^{-45},\qquad
\frac{\alpha_G}{\alpha}=2.4\times10^{-43}.
\end{equation}

\section{Falsifiability}
The model gives testable predictions. For magneto-gravitational devices
(Roshchin--Godin type) the relative correction to weight from precession is
estimated as
\begin{equation}
\frac{\Delta G}{G}\sim\frac{W R}{c^2}\sim 10^{-13},
\end{equation}
which is $\sim\!12$ orders of magnitude smaller than the claimed effects
($\sim\!35\%$); the latter, if real, would \emph{refute} the model, not confirm
it. This is a sign of a scientific, not a fitted, theory.

\section{Limits: proved / assumed / open}
\paragraph{Proved (with the derivation above).} The dimension of the charge
(Lemma~\ref{lem:qm}); the closure of mass; stability and annihilation
(Theorem~\ref{th:stab}); statistics from the period; the vanishing of the
current moment (Theorem~\ref{th:mu0}); the positivity of the gravity source;
the $1/R^2$ law (Theorem~\ref{th:invsq}); the necessity of two scales; the
Frank vector $\Omega=\pi$ of the twist (Lemma~\ref{lem:franck}).

\paragraph{Assumed.} The values $\alpha$, $\alpha_G$; $G$ (embedded in
$\gamma_P$); the two stiffness scales.

\paragraph{Open.} The numerical values of $\alpha,\alpha_G$; a rigorous $g=2$
(classically $g=1$; $g=2$ only at the spinor level \cite{osipov2010}); the
dynamical film stiffness (the dimensional factor in the source); the weak- and
strong-sector couplings.

\paragraph{The tensor ($J{=}2$) gravity channel --- a diagnosis.} The source
$\sigma\propto\langle W^2\rangle$ used above is the \emph{trace} (scalar, $J=0$)
of the quadrupole tensor $\langle W_iW_j\rangle$. The true tensor graviton
(spin-2 in Hehl's sense \cite{hehl1995}) lives in the trace-free part
$Q_{ij}=\langle W_iW_j\rangle-\tfrac13\delta_{ij}\langle W^2\rangle$. For a
single precessing ring with axis $\mathbf n$ a direct computation gives
$\langle W_iW_j\rangle=\tfrac12\omega^4a^2(\delta_{ij}-n_in_j)$, i.e.\
$Q_{ij}\neq0$ (an oblate quadrupole in the precession plane, eigenvalues
$\{-\tfrac13,+\tfrac16,+\tfrac16\}\omega^4a^2$). However, averaging over
isotropically oriented axes sends $\langle Q_{ij}\rangle\to0$: for an
unpolarized ensemble only the scalar $\langle W^2\rangle$ survives. Thus the
``scalarity'' of the model's gravity is a consequence of orientational
averaging, not of the absence of $J=2$. The tensor channel would manifest only
for \emph{polarized} (spin-aligned) matter. This is a diagnosis and a testable
prediction, not a derivation of the tensor graviton: the explicit field
equation for $Q_{ij}$ remains open.

\section{Conclusion}
A full derivation of all structural results of the model has been given. The
proofs are elementary (dimensional analysis, the Green function, the virial
theorem), but it was precisely their absence in earlier versions that led to
fitting. The proved, assumed and open parts are honestly separated. The next
steps are a rigorous spinor derivation of $g=2$ and this English version for
archival publication.

\appendix
\section{Numerical appendix: figures}
The computations are performed in model elastic units ($a=\omega=b=c=1$). The
absolute values are conventional; the physical content is in the \emph{decay
laws} and \emph{relations} proved in the main text.

\begin{figure}[h]
\centering
\includegraphics[width=0.75\linewidth]{figures_unified/channel_scaling.png}
\caption{Channel decay laws (log-log): Coulomb and gravity $\sim1/R^2$,
magnetic $\sim1/R^4$, exchange $\sim1/R^3$, weak --- exponential short range.
The slopes follow Theorem~\ref{th:invsq} and the dipole--dipole structure.}
\end{figure}

\begin{figure}[h]
\centering
\includegraphics[width=0.75\linewidth]{figures_unified/stability.png}
\caption{Stability of the pair (Theorem~\ref{th:stab}): $E(d)=-A/d+J^2/(2md^2)$.
The minimum at $d_{\text{eq}}=J^2/(mA)$ grows with the momentum $J$; as $J\to0$
the barrier disappears (annihilation).}
\end{figure}

\begin{figure}[h]
\centering
\includegraphics[width=0.75\linewidth]{figures_unified/twist_period.png}
\caption{Burgers twist: $\cos(\varphi/2)$ is antiperiodic under $2\pi$ and
returns only after $4\pi$ (spinor, fermion), unlike the ordinary $\cos\varphi$
with period $2\pi$ (boson).}
\end{figure}

\begin{thebibliography}{99}
\bibitem{fermi1923} E.~Fermi, Nuovo Cimento \textbf{25}, 159 (1923).
\bibitem{lazar2010} M.~Lazar, Math.\ Mech.\ Solids \textbf{16}(3), 253 (2010).
\bibitem{kadic1983} A.~Kadi\'c, D.~G.~B.~Edelen, Lecture Notes in Physics \textbf{174}, Springer (1983).
\bibitem{osipov2010} E.~A.~Kochetov, V.~A.~Osipov, arXiv:1001.4450 (2010).
\bibitem{hehl1995} F.~W.~Hehl, J.~D.~McCrea, E.~W.~Mielke, Y.~Ne'eman, Phys.\ Rep.\ \textbf{258}, 1 (1995); arXiv:gr-qc/9402012.
\bibitem{hirth1982} J.~P.~Hirth, J.~Lothe, \emph{Theory of Dislocations}, 2nd ed., Wiley (1982).
\bibitem{mistura1995} L.~Mistura, Int.\ J.\ Eng.\ Sci.\ \textbf{33}, 2149 (1995).
\bibitem{rohrlich1960} F.~Rohrlich, Am.\ J.\ Phys.\ \textbf{28}, 639 (1960).
\end{thebibliography}

\end{document}
